\documentclass[12pt]{article}

\usepackage{amsmath,amssymb,amsthm}
\usepackage{geometry}
\usepackage{setspace}
\usepackage{enumitem}
\usepackage{booktabs}
\usepackage{hyperref}

\geometry{letterpaper, margin=1in}
\setstretch{1.15}

\newtheorem{theorem}{Theorem}[section]
\newtheorem{proposition}[theorem]{Proposition}
\newtheorem{corollary}[theorem]{Corollary}
\newtheorem{lemma}[theorem]{Lemma}

\theoremstyle{definition}
\newtheorem{definition}[theorem]{Definition}

\theoremstyle{remark}
\newtheorem{remark}[theorem]{Remark}

\title{\textbf{The Arithmetic-Geometric Equivalence: \\[0.3cm]
Why the Cubic Lattice, the Lemniscate, \\
and the Prime Classification \\
Are One Structure}}

\author{W.\,J.\ cpaci-tani~III}
\date{}

\begin{document}

\maketitle

\begin{abstract}
\noindent We prove that four apparently independent mathematical objects---the cubic lattice $\mathbb{Z}^3$, the CM elliptic curve $E\!: y^2 = x^3 - x$, the Gaussian integer ring $\mathbb{Z}[i]$, and the Fermat two-square classification of primes---are locked together by a chain of theorems such that any one determines the other three. The chain passes through a single constant, $G^* = \sqrt{2}\,\Gamma(1/4)^2/(2\pi)$, which is simultaneously the Watson integral of the cubic lattice, the scaled period of the CM curve, and the theta function of the Gaussian lattice. We argue that if a physical theory's coupling constants are roots of a self-consistency equation involving $G^*$, then the ``axiom'' that space is $\mathbb{Z}^3$ is not a postulate but a \emph{tautology}: it is the geometric name for an arithmetic structure that exists independently.
\end{abstract}

\tableofcontents

%------------------------------------------------------------------
\section{Introduction: The Problem of the Axiom}
%------------------------------------------------------------------

Every physical theory begins with axioms. The Standard Model postulates gauge groups. General Relativity postulates a Lorentzian manifold. Quantum Mechanics postulates a Hilbert space. In each case, the axiom is accepted because the theory built on it matches experiment---but the axiom itself is not derived.

Other papers in this series derive the fine structure constant from a self-consistency equation on the cubic lattice $\mathbb{Z}^3$. The derivation requires stating ``space is $\mathbb{Z}^3$'' as an axiom. A hostile reviewer asks: \emph{why $\mathbb{Z}^3$? Why not a hexagonal lattice, or a BCC lattice, or a continuous manifold?}

This paper answers that question---not by deriving $\mathbb{Z}^3$ from something deeper, but by showing that $\mathbb{Z}^3$ is \emph{not an independent assumption}. It is one face of a four-way equivalence. The cubic lattice, the CM curve $E$, the Gaussian integers $\mathbb{Z}[i]$, and the Fermat prime classification are different descriptions of one mathematical object. Choosing any one determines the other three. The ``axiom'' is a tautology.

%------------------------------------------------------------------
\section{The Four Descriptions}
\label{sec:four}
%------------------------------------------------------------------

\begin{definition}[Description A: The Cubic Lattice]
$\mathbb{Z}^3$ is the set of integer-coordinate points in three-dimensional Euclidean space, with the standard nearest-neighbor structure (26-neighbor Moore neighborhood).
\end{definition}

\begin{definition}[Description B: The CM Curve]
$E\!: y^2 = x^3 - x$ is the unique elliptic curve over $\mathbb{Q}$ with $j$-invariant $j = 1728$, automorphism group $\operatorname{Aut}(E) \cong \mathbb{Z}/4\mathbb{Z}$, and complex multiplication by $\mathbb{Z}[i]$.
\end{definition}

\begin{definition}[Description C: The Gaussian Integers]
$\mathbb{Z}[i] = \{a + bi : a, b \in \mathbb{Z}\}$ is the ring of integers in the imaginary quadratic field $\mathbb{Q}(i)$, with norm $N(a+bi) = a^2 + b^2$.
\end{definition}

\begin{definition}[Description D: The Fermat Classification]
A prime $p$ is \emph{split} if $p \equiv 1 \pmod{4}$ (equivalently, $p = a^2 + b^2$ for some integers $a, b$), \emph{inert} if $p \equiv 3 \pmod{4}$, and \emph{ramified} if $p = 2$.
\end{definition}

%------------------------------------------------------------------
\section{The Equivalence Chain}
\label{sec:chain}
%------------------------------------------------------------------

We now prove that each description determines the next.

\subsection{A $\Rightarrow$ B: The lattice selects the curve}

\begin{theorem}[The cubic lattice forces the CM curve]
\label{thm:AtoB}
The coordinate planes of $\mathbb{Z}^3$ have $\mathbb{Z}_4$ rotational symmetry. Watson's evaluation of the lattice Green's function \cite{watson1939} reduces to complete elliptic integrals at a modulus $k$ determined by this symmetry. The unique elliptic modulus with $\mathbb{Z}_4$ automorphism symmetry is $k = 1/\sqrt{2}$ (the lemniscatic modulus), corresponding to the curve $E\!: y^2 = x^3 - x$ with $j = 1728$.
\end{theorem}

\subsection{B $\Rightarrow$ C: The curve determines the ring}

\begin{theorem}[The CM curve forces the Gaussian integers]
\label{thm:BtoC}
The endomorphism ring of $E$ over $\overline{\mathbb{Q}}$ is $\operatorname{End}(E) = \mathbb{Z}[i]$ \cite{silverman}. This is the maximal order in the imaginary quadratic field $\mathbb{Q}(i) = \mathbb{Q}(\sqrt{-1})$. The period lattice of $E$ in $\mathbb{C}$ is $\Lambda = \varpi(\mathbb{Z} + i\mathbb{Z})$, a square lattice with the same $\mathbb{Z}_4$ symmetry as the coordinate planes of $\mathbb{Z}^3$.
\end{theorem}

\subsection{C $\Rightarrow$ D: The ring classifies the primes}

\begin{theorem}[The Gaussian integers force the Fermat split]
\label{thm:CtoD}
A rational prime $p$ factors in $\mathbb{Z}[i]$ according to the Fermat--Euler criterion:
\begin{itemize}[nosep]
    \item $p \equiv 1 \pmod{4}$: $p$ \emph{splits} as $p = \pi \bar{\pi}$ where $\pi = a + bi$, $N(\pi) = p$.
    \item $p \equiv 3 \pmod{4}$: $p$ \emph{remains prime} (inert) in $\mathbb{Z}[i]$.
    \item $p = 2$: $p$ \emph{ramifies} as $2 = -i(1+i)^2$.
\end{itemize}
This is equivalent to Fermat's two-square theorem: $p = a^2 + b^2$ has a solution if and only if $p \not\equiv 3 \pmod{4}$.
\end{theorem}

\subsection{D $\Rightarrow$ A: The prime classification recovers the lattice}

\begin{theorem}[The Fermat split implies the cubic lattice]
\label{thm:DtoA}
The Fermat classification is the splitting behavior of primes in $\mathbb{Z}[i]$. The ring $\mathbb{Z}[i]$ is the endomorphism ring of the unique CM curve with $\mathbb{Z}_4$ symmetry ($j = 1728$). The Watson integral of a $D$-dimensional cubic lattice $\mathbb{Z}^D$ involves the Gamma function determined by the lattice's planar symmetry:
\begin{itemize}[nosep]
    \item $\mathbb{Z}^D$ with square cross-sections ($\mathbb{Z}_4$ symmetry) $\to$ $\Gamma(1/4)$ $\to$ $E\!: y^2 = x^3 - x$
    \item Hexagonal lattice ($\mathbb{Z}_6$ symmetry) $\to$ $\Gamma(1/3)$ $\to$ $E'\!: y^2 = x^3 - 1$ ($j = 0$)
\end{itemize}
Therefore: if the governing ring is $\mathbb{Z}[i]$ (which it must be, given the Fermat classification), the lattice must be cubic with $\mathbb{Z}_4$ planar symmetry. Among $D$-dimensional cubic lattices, $D = 3$ is the unique dimension where $\lfloor x_-(D) \rfloor = D$, closing the self-referential loop.
\end{theorem}

%------------------------------------------------------------------
\section{The Lock: $G^*$ at the Center}
\label{sec:lock}
%------------------------------------------------------------------

The four descriptions converge on a single constant.

\begin{theorem}[The Central Lock]
\label{thm:lock}
The constant $G^* = \sqrt{2}\,\Gamma(1/4)^2/(2\pi) \approx 2.9587$ is simultaneously:
\begin{enumerate}[nosep]
    \item The Watson integral of $\mathbb{Z}^3$: $G^{*2} = 2\pi W_3$ \quad (Description A)
    \item The scaled period of $E$: $G^* = 2\varpi/\sqrt{\pi}$ \quad (Description B)
    \item The theta function of $\mathbb{Z}[i]$: $G^* = \sqrt{2\pi}\;\theta_3(e^{-\pi})^2$ \quad (Description C)
    \item The generating function for the Fermat shell counts: $\theta_3(q)^2 = \sum r_2(n) q^n$ \quad (Description D)
\end{enumerate}
These four representations are proven to be equal. No circular reasoning is involved: each is derived from its own domain using independent methods (lattice Green's functions, elliptic integrals, modular forms, and arithmetic functions respectively), and the equality is a theorem.
\end{theorem}

\begin{corollary}[The Tautology]
\label{cor:tautology}
Postulating $\mathbb{Z}^3$ as the geometry of space is equivalent to postulating that $G^*$ is the fundamental coupling scale, which is equivalent to postulating $\mathbb{Z}[i]$ as the governing ring, which is equivalent to postulating the Fermat prime classification. These are four names for one mathematical object.
\end{corollary}

%------------------------------------------------------------------
\section{The Visible Evidence}
\label{sec:visible}
%------------------------------------------------------------------

The equivalence is not merely algebraic. It is \emph{visible} in the distribution of Gaussian primes.

\subsection{The mod-4 backbone}

The Gaussian prime condition (off-axis) requires $a^2 + b^2$ to be prime. Since $a^2 \equiv 0$ or $1 \pmod{4}$, the sum $a^2 + b^2$ can only be $0$, $1$, or $2 \pmod{4}$. For the sum to be an odd prime ($\equiv 1$ or $3 \pmod{4}$), exactly one of $a, b$ must be even and the other odd. This creates a period-4 checkerboard constraint on the lattice $\mathbb{Z}[i]$---the same $\mathbb{Z}_4$ symmetry that appears in Descriptions A--D.

\subsection{The Fermat channels}

The integers $6 = 2 \times 3$ and $7$ cannot be written as sums of two squares (both involve primes $\equiv 3 \pmod{4}$ to odd powers). Therefore $r_2(6) = r_2(7) = 0$, creating a \emph{representation gap} in the theta function $\theta_3(q)^2$. The cumulative count $R(5) = 20$ is the number of lattice points on $\mathbb{Z}^2$ within distance $\sqrt{5}$ of the origin---the ``knight's move'' shell. This gap is forced by Fermat's theorem and is visible in the moat structure of the Gaussian prime distribution as empty annular channels.

\subsection{The shell invariants}

The cumulative representation counts $R(N)$ hit the invariants of the cuboctahedron (the convex hull of the FCC sublattice of $\mathbb{Z}^3$):

\begin{center}
\renewcommand{\arraystretch}{1.2}
\begin{tabular}{@{}cccl@{}}
\toprule
$N$ & $R(N)$ & Lattice invariant & Description \\
\midrule
1 & 4 & $|\operatorname{Aut}(E)|$ & CM automorphism group \\
2 & 8 & $2^D$ & BCC corner count \\
4 & 12 & FCC neighbor count & Cuboctahedral vertices \\
5 & 20 & $b_3 + N_\mathrm{eff}$ & First Fermat plateau \\
8 & 24 & $|O|$ & Rotation group order \\
16 & 48 & $|O_h|$ & Full octahedral group \\
\bottomrule
\end{tabular}
\end{center}

These are not fitted. They are the cumulative lattice-point counts of $\mathbb{Z}^2$, which is the period lattice of $E$, which is the Gaussian integer ring $\mathbb{Z}[i]$, which classifies primes by the Fermat criterion. Each row is a theorem.

%------------------------------------------------------------------
\section{The Inert/Split Divide as Force Classification}
\label{sec:forces}
%------------------------------------------------------------------

The Fermat classification partitions the primes relevant to the master quadratic:

\begin{center}
\renewcommand{\arraystretch}{1.2}
\begin{tabular}{@{}cclll@{}}
\toprule
$p$ & Role & $p \bmod 4$ & $\mathbb{Z}[i]$ & $E \bmod p$ \\
\midrule
3 & $N_c$ (colors) & 3 & \textbf{inert} & supersingular \\
7 & $b_3$ (QCD beta) & 3 & \textbf{inert} & supersingular \\
47 & $D = |O_h| - 1$ & 3 & \textbf{inert} & supersingular \\
\midrule
13 & $N_\mathrm{eff}$ (DOF) & 1 & splits & ordinary \\
137 & $\approx 1/\alpha$ & 1 & splits & ordinary \\
\bottomrule
\end{tabular}
\end{center}

The integers associated with confinement $\{3, 7, 47\}$ are \emph{inert} primes in $\mathbb{Z}[i]$: the CM curve $E$ has supersingular reduction at these primes, and the Hecke eigenvalue vanishes ($a_p = 0$). The integers associated with electromagnetism $\{13, 137\}$ are \emph{split} primes: $E$ has ordinary reduction, and $137 = 4^2 + 11^2$ factors as $(4+11i)(4-11i)$ in $\mathbb{Z}[i]$.

This classification is a consequence of Fermat's theorem (1640, proved by Euler 1749). It predates FTD by four centuries. The observation that it separates the confinement integers from the electromagnetic integers is either a deep structural fact or a coincidence of extraordinary specificity.

%------------------------------------------------------------------
\section{Why This Is a Tautology}
\label{sec:tautology}
%------------------------------------------------------------------

The word ``tautology'' is usually pejorative. We use it precisely.

A tautology is a statement that is true by logical necessity---true in every possible interpretation. The equivalence A $\Leftrightarrow$ B $\Leftrightarrow$ C $\Leftrightarrow$ D is a tautology in this sense: given any one of the four descriptions, the other three follow by theorems that have been proven for 80--400 years (Fermat 1640, Gauss 1796, Watson 1939, Chowla--Selberg 1967).

The \emph{content} of a tautology is not that it is true---that is guaranteed---but that it \emph{reveals a connection that was not previously visible}. Euler's identity $e^{i\pi} + 1 = 0$ is a tautology (it follows from the definitions of $e$, $i$, $\pi$, $1$, and $0$). Its content is that exponentiation, complex numbers, and the circle are aspects of one structure.

Similarly: the content of the arithmetic-geometric equivalence is that the cubic lattice, the lemniscate, the Gaussian integers, and the prime classification are aspects of one structure. The master quadratic is a property of that structure. Its roots are properties of that structure. Whether those roots \emph{are} the coupling constants of our universe is a separate question---a physical conjecture, not a mathematical tautology.

But the tautology constrains the conjecture. If the coupling constants come from $G^*$, then space must be $\mathbb{Z}^3$, the primes must split by Fermat, the governing ring must be $\mathbb{Z}[i]$, and the CM curve must be $E$. All four follow, or none do. There is no room to accept the coupling constants while rejecting the lattice, or to accept the lattice while rejecting the prime classification. The structure is indivisible.

%------------------------------------------------------------------
\section{The Orthogonal Force Grid}
\label{sec:orthogonal}
%------------------------------------------------------------------

The two CM fields with non-trivial automorphisms---$\mathbb{Q}(i)$ and $\mathbb{Q}(\sqrt{-3})$---classify primes by two independent criteria: $p \bmod 4$ (splitting in $\mathbb{Z}[i]$) and $p \bmod 3$ (splitting in $\mathbb{Z}[\omega]$, $\omega = e^{2\pi i/3}$).

\begin{theorem}[Orthogonality]
\label{thm:orthogonal}
The mod-4 and mod-3 prime classifications are statistically independent. By the Chinese Remainder Theorem, $\gcd(3,4) = 1$ implies the joint distribution mod~12 is uniform over the four coprime residue classes $\{1, 5, 7, 11\}$.
\end{theorem}

\begin{proof}
Direct application of the Chinese Remainder Theorem and Dirichlet's theorem on primes in arithmetic progressions. \qedhere
\end{proof}

The four residue classes define an \emph{orthogonal force grid}:

\begin{center}
\renewcommand{\arraystretch}{1.2}
\begin{tabular}{@{}cccll@{}}
\toprule
Mod 12 & $\mathbb{Z}[i]$ & $\mathbb{Z}[\omega]$ & Sector & Framework integer \\
\midrule
$1$ & split & split & Both forces free & $13 = N_\mathrm{eff}$ \\
$5$ & split & inert & EM free, strong confined & $137 \approx 1/\alpha$ \\
$7$ & inert & split & EM confined, strong free & $7 = b_3$ \\
$11$ & inert & inert & Both forces confined & $47 = D$ \\
\bottomrule
\end{tabular}
\end{center}

The four framework primes $\{7, 13, 47, 137\}$ occupy all four quadrants of this grid---one integer per sector. This is verified by direct computation and confirmed by the uniform distribution of primes across the four classes.

\begin{remark}[The modulus 12 and the $j$-invariant]
The CRT modulus $12 = \mathrm{lcm}(3,4) = N_c \times N_\mathrm{base}$ is also the number of FCC neighbors (cuboctahedral vertices) and satisfies $12^3 = 1728 = j(E)$. The $j$-invariant of the CM curve that governs the entire framework is the cube of the modulus that classifies forces. Whether this is structural or coincidental is not resolved by this paper.
\end{remark}

%------------------------------------------------------------------
\section{The Irreducible Core: $\Gamma(1/4)$}
\label{sec:gamma}
%------------------------------------------------------------------

Beneath the four descriptions lies a single function evaluated at a single point.

\begin{theorem}[Everything from $\Gamma(1/4)$]
\label{thm:gamma}
Every constant in the framework is an algebraic expression in $\Gamma(1/4)$ and $\Gamma(1/2) = \sqrt{\pi}$:
\begin{align}
    \varpi &= \frac{\Gamma(1/4)^2}{2\sqrt{2}\,\Gamma(1/2)} \\[4pt]
    G^* &= \frac{\sqrt{2}\,\Gamma(1/4)^2}{2\,\Gamma(1/2)^2} \\[4pt]
    W_3 &= \frac{\Gamma(1/4)^4}{4\,\Gamma(1/2)^6} \\[4pt]
    \pi &= \Gamma(1/2)^2
\end{align}
The master quadratic $x^2 - 16G^{*2}x + 16G^{*3} = 0$ and its roots $x_+ = 137.036\ldots$, $x_- = 3.024\ldots$ are therefore functions of $\Gamma(1/4)$ alone \emph{(}with $\pi = \Gamma(1/2)^2$ providing the circular period\emph{)}.
\end{theorem}

\begin{proof}
Direct substitution from the definitions in Section~\ref{sec:four} and the identity $\Gamma(1/2) = \sqrt{\pi}$ (Euler's reflection formula at $z = 1/2$). \qedhere
\end{proof}

\begin{remark}[Why $\Gamma$ and why $1/4$]
The Gamma function is not a choice. By the Bohr--Mollerup theorem (1922), $\Gamma$ is the \emph{unique} function on $(0,\infty)$ satisfying: (i)~$\Gamma(1) = 1$, (ii)~$\Gamma(x+1) = x\,\Gamma(x)$, (iii)~$\log\Gamma$ is convex. No other function has all three properties. The evaluation point $1/4$ is selected by the $\mathbb{Z}_4$ symmetry of the cubic lattice: the lemniscatic modulus $k = 1/\sqrt{2}$ forces the complete elliptic integral $K(1/\sqrt{2}) = \Gamma(1/4)^2/(4\sqrt{\pi})$, and the Chowla--Selberg formula evaluates the theta function at $\tau = i$ in terms of $\Gamma$ at arguments determined by the CM discriminant $-4$, giving $\Gamma(1/4)$.

The argument $1/3$ would give $\Gamma(1/3)$, the hexagonal lattice ($\mathbb{Z}_6$ symmetry), and the CM curve $y^2 = x^3 - 1$ with $j = 0$. The argument $1/4$ gives $\Gamma(1/4)$, the cubic lattice ($\mathbb{Z}_4$ symmetry), and the CM curve $y^2 = x^3 - x$ with $j = 1728$. These are the only two imaginary quadratic fields of class number~1 with non-trivial automorphisms ($\mathbb{Q}(i)$ and $\mathbb{Q}(\sqrt{-3})$). The cubic lattice corresponds to $1/4$; the hexagonal to $1/3$. There is no third option.
\end{remark}

\begin{corollary}[The argument from nothing]
The chain from no assumptions to the master quadratic roots is:
\begin{equation}
    \underset{\text{Bohr--Mollerup}}{\Gamma} \;\xrightarrow{\;\mathbb{Z}_4\;}\; \Gamma(1/4) \;\xrightarrow{\;\text{algebra}\;}\; G^* \;\xrightarrow{\;\text{self-consistency}\;}\; x_+ = 137.036
\end{equation}
The first arrow is a theorem (the unique log-convex factorial extension). The second is a theorem (Watson 1939, Chowla--Selberg 1967). The third is algebra. The fourth is the self-consistency prescription \emph{[SELECTION]}---the one remaining non-tautological step in the entire framework.
\end{corollary}

%------------------------------------------------------------------
\section{Summary}
%------------------------------------------------------------------

\begin{enumerate}
    \item Four mathematical objects---$\mathbb{Z}^3$, $E\!: y^2 = x^3 - x$, $\mathbb{Z}[i]$, and the Fermat classification---are equivalent descriptions of one structure, locked together by theorems spanning four centuries.

    \item The constant $G^*$ sits at the center: it is the Watson integral of $\mathbb{Z}^3$, the period of $E$, the theta function of $\mathbb{Z}[i]$, and the generating function for the Fermat shell counts.

    \item The ``axiom'' that space is $\mathbb{Z}^3$ is not an independent postulate. It is the geometric name for the arithmetic structure that the primes already exhibit. Accepting it is no more (and no less) than accepting that $\mathbb{Z}[i]$ is the governing ring.

    \item Whether the physical universe lives in this structure is a conjecture. But the structure itself is a tautology---true by mathematical necessity, visible in the prime distribution, and indivisible.
\end{enumerate}

%------------------------------------------------------------------
\begin{thebibliography}{9}
%------------------------------------------------------------------

\bibitem{watson1939}
G.\,N.\ Watson, ``Three triple integrals,'' \textit{Quart.\ J.\ Math.\ Oxford} \textbf{10} (1939), 266--276.

\bibitem{silverman}
J.\,H.\ Silverman, \textit{The Arithmetic of Elliptic Curves}, 2nd ed., Springer GTM~\textbf{106} (2009).

\bibitem{chudnovsky}
G.\,V.\ Chudnovsky, ``Contributions to the theory of transcendental numbers,'' \textit{Math.\ Surveys and Monographs} \textbf{19}, AMS (1984).

\end{thebibliography}

\end{document}
